\documentclass[multi=page, border=15pt, varwidth=16cm]{standalone}
\usepackage{ctex}      % 提供中文支持
\usepackage{tabularx}  % 自动调整列宽
\usepackage{booktabs}  % 优雅的三线表支持
\usepackage{xcolor}
\usepackage{colortbl}

% 统一的深灰蓝表头色
\definecolor{headerbg}{RGB}{52, 73, 94}
\definecolor{headertext}{RGB}{255, 255, 255}

% 自定义表头文本样式命令
\newcommand{\thnode}[1]{\textcolor{headertext}{\bfseries #1}}

\begin{document}
\sffamily
\renewcommand{\arraystretch}{1.5}

% ================= 图片：fsck 修复工具集 =================
\begin{page}
	\centering
	% \textbf{\Large fsck 文件系统检查与修复参数} \\ \vspace{0.8em}

	\begin{tabularx}{\linewidth}{>{\bfseries\color{blue!70!black}}c X >{\bfseries\color{blue!70!black}}c X >{\bfseries\color{blue!70!black}}c X}
		\toprule
		\rowcolor{headerbg}
		\thnode{参数} & \thnode{说明}          & \thnode{参数} & \thnode{说明}      & \thnode{参数} & \thnode{说明} \\
		\midrule
		% 第一行：核心修复
		-y          & \textbf{自动修复错误}      & -a          & 自动修复(同-y)        & -r          & 出现错误时提示     \\
		% 第二行：批量与自动化
		-A          & \textbf{检查fstab全部分区} & -R          & 跳过根文件系统检查        & -s          & 依次检查多个分区    \\
		% 第三行：显示与反馈
		-V          & 显示详细输出信息             & -C          & \textbf{显示修复进度条} & -T          & 不显示头部标题     \\
		% 第四行：测试与特殊
		-N          & 仅模拟执行不改磁盘            & -t          & 指定文件系统类型         & -M          & 不检查已挂载设备    \\
		\bottomrule
		\multicolumn{6}{p{\linewidth}}{\small \color{gray}\textit{* 注：\texttt{-y} 是最常用的参数，能避免在修复成千上万个损坏节点时手动输入 yes}}
	\end{tabularx}

	% \vspace{1.5em}

	% 安全警告栏
	\begin{tabularx}{\linewidth}{>{\columncolor{red!10}}l X}
		\toprule
		% \rowcolor{red!60!black}
		% \thnode{核心禁忌} & \thnode{绝对禁止在已挂载（Mounted）的文件系统上运行 \texttt{fsck}，否则会导致数据毁灭性损坏！}                                                     \\
		% \midrule
		\textbf{正确步骤} & \texttt{sudo umount /dev/sdb1} $\rightarrow$ \texttt{sudo fsck -y /dev/sdb1} $\rightarrow$ \texttt{sudo mount ...} \\
		\bottomrule
	\end{tabularx}
\end{page}

\end{document}